\documentclass{standalone}
\usepackage{pgfplots}

\usetikzlibrary{positioning}

\pgfplotsset{compat=1.18}
\pgfplotsset{plot coordinates/math parser=false,unbounded coords=discard}

\providecommand{\dataDir}{paper_main}
% \providecommand{\zoomSubdirs}{94-108_9-25}
% \providecommand{\zoomSubdirs}{-8-9_-6-6,94-108_9-25}
% \providecommand{\zoomSubdirs}{94-108_9-25,-8-9_-6-6}

\begin{document}
\begin{tikzpicture}
  \begin{axis}[
    xlabel={t-SNE${}_1$},
    ylabel={t-SNE${}_2$},
    title={Latent space projection (t-SNE)},
    grid=major,
    clip=true,
    scatter/classes={
      Unlabeled={
        mark=o,
        draw=gray,
        draw opacity=0.05,
        fill=none,
      },
      {NO Event}={
        mark=diamond*,
        fill=red,
      },
      Event={
        mark=square*,
        fill=green
      }
    },
  legend style={
    at={(0.5,-0.2)}, % adjust position
    anchor=north,     % attach legend’s top to that point
    legend columns=-1 % put all entries in one horizontal row
  },
  ]

    % one pass: use the "Color" column to pick class styles
    \addplot[
      scatter,
      only marks,
      scatter src=explicit symbolic,
    ]
    table[
      x=tsne_1, y=tsne_2,
      meta=Color,
      col sep=comma,
      trim cells=true,
      ignore chars={ },
    % ]{\dataDir/foo.csv};
    ]{\dataDir/tSNE.csv};

    % ---------------- Legend entries ----------------
    % \addlegendimage{}
    % \addlegendentry{}

    \addlegendimage{only marks, /pgfplots/scatter/classes/Unlabeled}
    % \addlegendimage{
    %   only marks, /pgfplots/scatter/classes/Unlabeled
    %   fill=black,
    % }
    \addlegendentry{Unlabeled}

    \addlegendimage{only marks, /pgfplots/scatter/classes/{NO Event}} % note braces for space
    \addlegendentry{No Event}

    \addlegendimage{only marks, /pgfplots/scatter/classes/Event}
    \addlegendentry{Event}

    % if zoom subdirs (comma-separated list like "-8-9_-6-6,-95--84_58-66") have been passed...  
    \ifdefined\zoomSubdirs

      % for each subdir...
      \foreach \subdir [count=\iSubdir from 1] in \zoomSubdirs {
        
        % CSV with the coordinates of the zoom rectangle are read
        \pgfplotstableread[col sep=comma]{\dataDir/\subdir/coordinates.csv}\zoombox

        % corners are parsed from there
        \pgfplotstablegetelem{0}{x}\of\zoombox \edef\xa{\pgfplotsretval}
        \pgfplotstablegetelem{0}{y}\of\zoombox \edef\ya{\pgfplotsretval}
        \pgfplotstablegetelem{1}{x}\of\zoombox \edef\xb{\pgfplotsretval}
        \pgfplotstablegetelem{1}{y}\of\zoombox \edef\yb{\pgfplotsretval}

        % \typeout{[zoom \iSubdir] \subdir: xa=\xa, ya=\ya, xb=\xb, yb=\yb}

        % rectangle is plotted...
        \expanded{\noexpand\draw[black, thick] (axis cs:\xa,\ya) rectangle (axis cs:\xb,\yb);}

        % ...along with a "label"
        \expanded{\noexpand\node[right] at (axis cs:\xa,\ya) {\#\iSubdir};}

        % highlighted points for this region are plotted
        \addplot+[
          only marks,
          mark=star,
          mark size=2.2pt,
          draw=blue,
        ] table[
          x=tsne_1, y=tsne_2,
          col sep=comma,
        ]{\dataDir/\subdir/selected_windows.csv};
      }

    \fi

  \end{axis}
\end{tikzpicture}

\end{document}
